\section{NoSQL and Distributed Databases}

% ----------------------------------------------------------
\begin{frame}{ACID -- The Gold Standard for Transactions}
  \begin{block}{ACID Properties}
    The four guarantees that define a \textbf{reliable database transaction} in classical
    relational databases:
  \end{block}

  \begin{itemize}
    \item \textbf{Atomicity}: a transaction is \emph{all-or-nothing}. Either every operation
          succeeds and is committed, or none of them are applied.
          Implemented via before-images and \texttt{ROLLBACK}.

    \item \textbf{Consistency}: a transaction always moves the database from one \emph{valid state}
          to another. All constraints (PKs, FKs, CHECK) must hold before and after.

    \item \textbf{Isolation}: concurrent transactions behave \emph{as if executed serially}.
          In-progress state of one transaction is invisible to others.
          Implemented via locks or MVCC (Multi-Version Concurrency Control).

    \item \textbf{Durability}: once a transaction commits, its changes are \emph{permanent} --
          even if the system crashes immediately afterward.
          Implemented via the write-ahead transaction log (WAL) and checkpoints.
  \end{itemize}

  \begin{exampleblock}{ACID Systems}
    MySQL, PostgreSQL, Oracle, SQL Server -- fully ACID-compliant.
    Default choice for financial, healthcare, and any domain where correctness is non-negotiable.
  \end{exampleblock}
\end{frame}

% ----------------------------------------------------------
\begin{frame}{CAP Theorem -- The Three Properties}
  \begin{block}{CAP Theorem (Eric Brewer, 2000)}
    Any \textbf{distributed} data store can simultaneously guarantee \textbf{at most two} of the
    following three properties:
  \end{block}

  \vspace{6pt}
  \textbf{C -- Consistency}
  \begin{itemize}
    \item Every read receives the \emph{most recent write}, or an error.
    \item All nodes in the system see the same data at the same time.
    \item Example: after a write to node A, node B immediately returns that value too.
  \end{itemize}

  \vspace{4pt}
  \begin{block}{C is about correctness of reads -- not the same as ACID Consistency.}
  \end{block}
\end{frame}

% ----------------------------------------------------------
\begin{frame}{CAP Theorem -- Availability and Partition Tolerance}
  \textbf{A -- Availability}
  \begin{itemize}
    \item Every request sent to a non-failing node receives a \emph{response} -- not necessarily
          the most recent data, but \emph{some} valid response (no timeouts, no errors).
    \item The system remains operational under all circumstances (as long as the node itself is up).
  \end{itemize}

  \vspace{8pt}
  \textbf{P -- Partition Tolerance}
  \begin{itemize}
    \item The system continues to operate correctly even if network messages between nodes
          are lost, delayed, or the network splits into isolated segments.
  \end{itemize}

  \begin{alertblock}{Partitions Will Happen}
    In any real distributed system (multiple machines, data centres), network partitions are
    \textbf{inevitable} -- cables fail, switches reboot, data centres lose connectivity.
    Therefore \textbf{P is not optional}. The real trade-off is always \textbf{C vs A}.
  \end{alertblock}
\end{frame}

% ----------------------------------------------------------
\begin{frame}{CAP Theorem -- CP vs AP Trade-off}
  When a network partition occurs, the system must choose:

  \begin{block}{CP -- Sacrifice Availability, Maintain Consistency}
    Cancel or reject operations that cannot be guaranteed consistent across all nodes.
    The system returns an \textbf{error} rather than a potentially stale answer.
    \begin{itemize}
      \item \textbf{Example}: MongoDB (default config) refuses writes to a
            partition-isolated node. HBase, Zookeeper.
      \item \textbf{Use case}: financial transactions, bookings -- stale data is unacceptable.
    \end{itemize}
  \end{block}

  \begin{block}{AP -- Sacrifice Consistency, Maintain Availability}
    Continue serving requests from all nodes even if some return stale data.
    The system will \emph{eventually converge} once the partition heals.
    \begin{itemize}
      \item \textbf{Example}: Cassandra, CouchDB, DynamoDB continue serving all nodes.
      \item \textbf{Use case}: social media feeds, product catalogs -- brief stale data is acceptable.
    \end{itemize}
  \end{block}
\end{frame}

% ----------------------------------------------------------
\begin{frame}{CAP Theorem -- System Classification}
  \begin{center}
  \renewcommand{\arraystretch}{1.35}
  \begin{tabularx}{\linewidth}{@{}llXl@{}}
    \toprule
    \textbf{Category} & \textbf{Examples} & \textbf{Trade-off} & \textbf{Partition?} \\
    \midrule
    CA & MySQL, PostgreSQL, Oracle & Consistent + Available & Single node only \\
    CP & MongoDB, HBase, Zookeeper & Consistent + Part.-tolerant & May reject requests \\
    AP & Cassandra, CouchDB, DynamoDB & Available + Part.-tolerant & May serve stale data \\
    \bottomrule
  \end{tabularx}
  \end{center}

  \begin{alertblock}{CA is Not Truly Distributed}
    CA systems (MySQL, PostgreSQL) achieve both C and A only because they run on a single node
    -- if there is no partition to tolerate, the theorem does not apply. Multi-node deployments
    of these systems (MySQL replication) introduce CA trade-offs.
  \end{alertblock}

  \begin{block}{Modern Nuance}
    Real systems are not binary. Many newer systems (Google Spanner, CockroachDB) offer
    ``essentially CP'' with tunable consistency levels, blurring the categories.
  \end{block}
\end{frame}

% ----------------------------------------------------------
\begin{frame}{What is NoSQL?}
  \begin{block}{Definition}
    \textbf{NoSQL} (``Not only SQL'') refers to database management systems designed for
    use cases not well-served by traditional relational databases. They typically trade
    strict ACID compliance for horizontal scalability, schema flexibility, or specialised
    data models.
  \end{block}

  \textbf{Why NoSQL emerged:}
  \begin{itemize}
    \item \textbf{Horizontal scalability}: relational DBs are hard to shard across hundreds of
          commodity servers. NoSQL DBs are designed to distribute from the ground up.
    \item \textbf{Schema flexibility}: document stores accept any JSON structure; no
          \texttt{ALTER TABLE} migration required for schema changes.
    \item \textbf{High write throughput}: logging millions of IoT events per second would
          overwhelm a single relational DB; distributed NoSQL clusters handle this natively.
    \item \textbf{Specialised data models}: graphs, documents, and key-value pairs map poorly
          to flat relational tables.
  \end{itemize}
\end{frame}

% ----------------------------------------------------------
\begin{frame}{NoSQL -- Data Model Comparison}
  \begin{center}
  \renewcommand{\arraystretch}{1.3}
  \begin{tabularx}{\linewidth}{@{}lXXl@{}}
    \toprule
    \textbf{Model} & \textbf{Data Structure} & \textbf{Use Case} & \textbf{Example} \\
    \midrule
    Document Store
      & JSON / BSON documents; flexible schema
      & Content management, user profiles, product catalogs
      & MongoDB, CouchDB \\[4pt]
    Key-Value Store
      & Hash map: key $\to$ opaque value; ultra-fast
      & Caching, session storage, leaderboards
      & Redis, DynamoDB \\[4pt]
    Columnar (Wide Column)
      & Rows with many sparse columns; column families
      & Time-series, IoT, analytics, write-heavy workloads
      & Cassandra, HBase \\[4pt]
    Graph Database
      & Nodes and edges with properties; optimised traversal
      & Social networks, recommendation engines, fraud detection
      & Neo4j, Amazon Neptune \\
    \bottomrule
  \end{tabularx}
  \end{center}
\end{frame}

% ----------------------------------------------------------
\begin{frame}{BASE -- The NoSQL Alternative to ACID}
  \begin{block}{BASE Properties}
    Most NoSQL databases follow the \textbf{BASE} model. BASE deliberately relaxes consistency
    guarantees in exchange for availability and performance:
  \end{block}

  \begin{itemize}
    \item \textbf{B}asically \textbf{A}vailable: the system is always available to accept
          requests. All users can access data concurrently without waiting for locks.
          The system prioritises responding over guaranteeing freshness.

    \item \textbf{S}oft State: the state of the system may change \emph{over time even without
          new user input}, as replication propagates across nodes in the background.
          Different nodes may temporarily hold different versions of the same record.

    \item \textbf{E}ventually Consistent: after all updates have stopped and sufficient
          propagation time has passed, all replicas \emph{will converge} to the same value.
          The system does NOT guarantee identical data on all nodes at the same instant.
  \end{itemize}

  \begin{exampleblock}{When is BASE Acceptable?}
    Social media timelines, product catalogs, shopping carts, recommendation engines --
    brief inconsistency is tolerable. Bank balances, medical records, payments -- ACID required.
  \end{exampleblock}
\end{frame}

% ----------------------------------------------------------
\begin{frame}{ACID vs BASE -- Comparison}
  \begin{center}
  \renewcommand{\arraystretch}{1.3}
  \begin{tabularx}{\linewidth}{@{}lXX@{}}
    \toprule
    \textbf{Property} & \textbf{ACID} & \textbf{BASE} \\
    \midrule
    Primary focus   & Data correctness and integrity & Availability and performance \\
    Consistency     & Guaranteed immediately after every commit & Eventually consistent (may lag) \\
    Transactions    & Full ACID with rollback & None, or limited / per-document \\
    Intermediate state & Always consistent & Soft state; temporary inconsistency expected \\
    Scalability     & Vertical (scale up: more CPU/RAM) & Horizontal (scale out: more servers) \\
    Use cases       & Finance, ERP, healthcare, banking & Social media, IoT, event logging, caches \\
    Examples        & MySQL, PostgreSQL, Oracle & MongoDB, Cassandra, Redis, DynamoDB \\
    \bottomrule
  \end{tabularx}
  \end{center}

  \begin{block}{Practical Takeaway}
    These are not competing philosophies -- they are different tools for different problems.
    Modern architectures often use \textbf{both}: a relational ACID DB for transactional data
    and a NoSQL BASE store for high-throughput events or caching.
  \end{block}
\end{frame}

% ----------------------------------------------------------
\begin{frame}{Real Systems and Their CAP / BASE Position}
  \begin{center}
  \renewcommand{\arraystretch}{1.35}
  \begin{tabularx}{\linewidth}{@{}llXX@{}}
    \toprule
    \textbf{System} & \textbf{Type} & \textbf{CAP Position} & \textbf{Key Strength} \\
    \midrule
    MySQL / MariaDB & Relational & CA (single node) & Full ACID, SQL \\
    PostgreSQL & Relational & CA (single node) & Advanced SQL, extensions \\
    MongoDB & Document & CP (tunable) & Flexible schema, rich queries \\
    Cassandra & Columnar & AP & Massive write throughput, HA \\
    Redis & Key-Value & AP (in-memory) & Sub-millisecond reads, caching \\
    DynamoDB & Key-Value & AP & Serverless, auto-scaling \\
    CouchDB & Document & AP & Offline-first, sync \\
    Neo4j & Graph & CA & Graph traversal performance \\
    \bottomrule
  \end{tabularx}
  \end{center}

  \begin{block}{No Single ``Best'' Database}
    Choose based on your data model, consistency requirements, scalability needs,
    and operational complexity budget.
  \end{block}
\end{frame}
